\documentclass[11pt,a4paper]{article}

\usepackage[T1]{fontenc}
\usepackage[utf8]{inputenc}
\usepackage[english]{babel}
\usepackage[a4paper,margin=2.2cm]{geometry}
\usepackage{lmodern}
\usepackage{microtype}
\usepackage{enumitem}
\usepackage{xcolor}
\usepackage{hyperref}
\usepackage{titlesec}
\usepackage{parskip}

\hypersetup{
  colorlinks=true,
  linkcolor=blue!55!black,
  urlcolor=blue!65!black,
  pdftitle={DIY e-Handpan Assembly Instructions},
  pdfsubject={Assembly manual},
  pdfauthor={}
}

\setlist[itemize]{leftmargin=1.6em,itemsep=0.25em,topsep=0.4em}
\setlist[enumerate]{leftmargin=1.9em,itemsep=0.3em,topsep=0.4em}
\titleformat{\section}{\Large\bfseries}{\thesection.}{0.6em}{}
\titleformat{\subsection}{\large\bfseries}{}{0em}{}
\titleformat{\subsubsection}{\normalsize\bfseries}{}{0em}{}

\newcommand{\tutorial}[1]{\textbf{Video tutorial:} \url{#1}}
\newcommand{\tools}{\subsection{Required tools}}
\newcommand{\materials}{\subsection{Materials}}
\newcommand{\procedure}{\subsection{Procedure}}
\newcommand{\note}[1]{\begin{quote}\textbf{Note:} #1\end{quote}}
\newcommand{\alternative}[1]{\begin{quote}\textbf{Alternative:} #1\end{quote}}

\title{DIY e-Handpan Assembly Instructions}
\author{Benoit Collin, Dominique Fourer and Eric Genotelle}
\date{\today}

\begin{document}
\maketitle
\tableofcontents
\newpage

\section*{Workshop Resources}
\addcontentsline{toc}{section}{Workshop Resources}

\textbf{Summary table of the assembly stages, required tools and materials, and estimated completion time for a group of 12 students:}

\url{https://docs.google.com/spreadsheets/d/16BShQwQICNdTbGfFbWgz5b8PJe-1twqu/edit?usp=drive_link&ouid=109005322539614617821&rtpof=true&sd=true}

\textbf{Playlist containing all video tutorials:}

\url{https://youtube.com/playlist?list=PL_k3PczvzYDb-f5t7U6dC4wk6cq4hSw6F&si=5ZHYeenZAPfkI9HT}

\section{Drilling the Bowls}
\tutorial{https://youtu.be/Srb9kXRlCo4}

\tools
\begin{itemize}
  \item A3 printer
  \item Scissors
  \item Permanent marker suitable for smooth surfaces
  \item Drill
  \item 3--5 mm drill bit
\end{itemize}

\materials
\begin{itemize}
  \item One or two A3 sheets
  \item 40 cm and 45 cm bowls
  \item Single-sided adhesive tape
\end{itemize}

\procedure
\subsubsection{Preparing the templates}
Print the templates on A3 paper. Double-sided printing may be used. Cut along the circular outline, then make a hole in the centre of each template.

\subsubsection{Marking the bowls}
\begin{itemize}
  \item Place the bowls upside down on a table.
  \item Drill a hole in the centre of each bowl.
  \item Align the centre of each template with the centre of its bowl by passing the drill bit through the template's central hole and then into the bowl's central hole.
\end{itemize}

\note{Use the ABC template for the 45 cm bowl and the B$\flat$--D$\flat$--E$\flat$ template for the 40 cm bowl.}

\begin{itemize}
  \item Secure each template to its bowl at two or three points with single-sided adhesive tape. This prevents the template from rotating during marking and is very important.
  \item Mark the bowls 8 cm from the rim and write the note name beside each hole on the outside of the bowls.
  \item Drill the bowls at the marked positions. Then write the note name beside each hole on the inside, ensuring that it matches the name written outside.
\end{itemize}

\section{Cutting and Piercing the Foam Discs}
\tutorial{https://youtu.be/3l_8JzbdJwk}

\tools
\begin{itemize}
  \item Scissors
  \item CD
  \item Felt-tip pen
\end{itemize}

\materials
\begin{itemize}
  \item Enough 1 cm-thick foam to cut twelve discs, each 12 cm in diameter
  \item Enough 2 cm-thick foam to cut two discs, each 12 cm in diameter
\end{itemize}

\procedure
\subsubsection{Drawing the circles on the foam}
Use the CD as a tracing template. Position it near the edge of the foam sheet, leaving an overhang of approximately 0.5--1 cm, and trace around the CD.

\subsubsection{Cutting and piercing the discs}
\begin{itemize}
  \item Cut inside the traced line, leaving a margin of approximately 0.5--1 cm.
  \item Make a hole in each disc approximately 1 cm from its centre, using the tip of the scissors.
\end{itemize}

\section{Attaching the Sensor and Foam Disc to the CD}
\tutorial{https://youtu.be/R5TMR4DGeH0?feature=shared}

\tools
\begin{itemize}
  \item Scissors
  \item A rigid drinking-straw-type tube, made of cardboard or plastic, with a maximum diameter of 5 mm
\end{itemize}

\materials
\begin{itemize}
  \item Piezoelectric sensors with pre-soldered wires
  \item Pierced foam discs
  \item 5 cm-wide double-sided adhesive tape
\end{itemize}

\procedure
\subsubsection{Making the ``straw'' tool}
Cut one end of the straw at an angle. The resulting piece should be the same length as the sensor wires.

\subsubsection{Applying adhesive to the CD}
\begin{itemize}
  \item Cut two trapezoids from the double-sided adhesive tape. Each piece must be smaller than the CD, with approximate dimensions $L=13$ cm, $l=9$ cm and $h=5$ cm.
  \item Apply one trapezoid to the CD, aligning the long edge of the strip with the CD's diameter so that it crosses the centre hole.
  \item Apply the second trapezoid beside the first so that the centre hole is covered. Do not overlap the two adhesive pieces.
  \item Trim any adhesive extending beyond the edge of the CD.
\end{itemize}

\subsubsection{Attaching the sensor and foam disc}
\begin{itemize}
  \item Remove the protective films from the adhesive trapezoids attached to the CD.
  \item Position the sensor at the centre of the CD and press lightly to bond them together.
  \item Feed the sensor wires through the straw, with the angled end of the straw at the free ends of the wires and the straight end facing the sensor. Twist the wires together to make insertion easier.
  \item Pass the straw through the hole in the foam disc.
  \item Slide the foam disc towards the CD while keeping it centred, using the straw as a guide. The hole in the disc must align with the solder joints where the wires meet the sensor.
  \item Check that the foam disc is centred on the CD, then press the foam firmly onto the CD to bond them together.
\end{itemize}

\section{Attaching the Striking Surfaces (Pads) to the Bowls}
\tutorial{https://youtu.be/PlesHf17iok?feature=shared}

\tools
\begin{itemize}
  \item Scissors
  \item Cardboard to protect the work surface
\end{itemize}

\materials
\begin{itemize}
  \item Spray adhesive or double-sided adhesive tape
  \item Adhesive and expanding-foam remover, such as acetone
\end{itemize}

\note{Attach the pads in the following order: begin with the central note, using the pad made with 2 cm-thick foam. Then attach the peripheral pad bearing the same note name, which is the octave, using a 1 cm-thick foam pad. Before proceeding, draw a point on the edge of the central pad, directly opposite and aligned with the octave pad.}

\procedure
\begin{itemize}
  \item Twist the wires together to make it easier to feed them through the hole in the bowl.
  \item Cut a 2 cm piece of double-sided adhesive tape. Apply it to the pad's foam as close as possible to the wires, or apply it to the bowl as close as possible to a hole. Remove the protective film.
  \item Place the bowl upside down on a table.
  \item Feed the wires through the hole. The point where the wires exit the pad must face upwards.
  \item Position the pad on the bowl and press firmly.
\end{itemize}

\alternative{When using spray adhesive, spray the adhesive around the corresponding hole in the bowl instead of using double-sided tape.}

\section{Making and Attaching the Supports}
\tutorial{https://youtu.be/RaQdDFHfXVI?feature=shared}

\tools
\begin{itemize}
  \item Scissors
  \item Permanent marker
  \item Cardboard to protect the work surface
\end{itemize}

\materials
\begin{itemize}
  \item Spray adhesive or double-sided adhesive tape
  \item Foam offcuts or expanding foam
  \item Adhesive and expanding-foam remover, such as acetone
\end{itemize}

\procedure
\begin{itemize}
  \item Cut crescent-shaped pieces from the foam offcuts.
  \item Cut a 2 cm piece of double-sided adhesive tape, apply it to a foam crescent, and remove the protective film.
  \item Insert the foam crescent beneath the pad, with its pointed ends facing the centre of the bowl. Press it firmly against the bowl.
\end{itemize}

\alternative{When using spray adhesive, spray adhesive onto the bowl beneath the pad at the point where the pad is furthest from the bowl, approximately 5 cm from the rim. This method is stronger but requires drying time and greater care when handling the bowl during construction. It also makes disassembly more difficult if a repair is needed.}

\alternative{The foam crescents may be replaced with expanding foam. This gives the pads better adhesion and stability, but requires additional drying time. Any excess foam must be trimmed after it has fully dried.}

\section{Drilling the Bowls for the LEDs}
\tutorial{https://youtu.be/Eq6pYfNYuWM?feature=shared}

\tools
\begin{itemize}
  \item Drill
  \item Drill bit matching the diameter of the LEDs
  \item Cloth or vacuum cleaner
\end{itemize}

\materials
No additional materials are required.

\procedure
\begin{itemize}
  \item Place the bowl upside down, with its rim resting on the table.
  \item Drill one hole near the edge of each CD, less than 1 cm away, on the side facing the centre of the bowl.
  \item Remove the drilling dust with a cloth or vacuum cleaner, then carefully remove any splinters or burrs from the holes.
\end{itemize}

\section{Inserting the LEDs into the Bowls}
\tutorial{https://youtu.be/NGQcG9u3N-Y?is=k9aHxWO3TmvNvXqQ}

\tools
\begin{itemize}
  \item Side cutters
\end{itemize}

\materials
\begin{itemize}
  \item Transparent instant adhesive
  \item Adhesive remover
\end{itemize}

\procedure
\begin{itemize}
  \item Shorten the shorter lead of each LED by 2--3 mm so that it can be identified easily.
  \item Apply adhesive to the LED and to the hole.
  \item Insert the LED into the hole with the short lead on the left.
\end{itemize}

\section{Ground Circuit}
\tutorial{https://youtu.be/tKWcBon72ZY?is=V0iVByUxMFY0HdpO}

\tools
\begin{itemize}
  \item Side cutters
  \item Soldering iron
  \item Damp sponge
  \item Wire stripper
  \item Scissors
  \item Utility knife
  \item Permanent marker
  \item Measuring tape
\end{itemize}

\materials
\begin{itemize}
  \item Solder
  \item Heat-shrink tubing
  \item Ground wire
\end{itemize}

\procedure
\subsubsection{Connecting the sensor grounds to the LEDs}
\begin{itemize}
  \item Cut fourteen 2 cm lengths of heat-shrink tubing.
  \item Strip the sensor ground wires if necessary.
  \item Bend the centre lead, or cathode, of each LED towards the ground wire of the corresponding sensor.
  \item Pass the central sensor's ground wire through one piece of heat-shrink tubing.
  \item Hold the end of the wire against the LED's centre lead, using a tool if necessary, and solder the joint.
  \item Repeat for every note.
\end{itemize}

\subsubsection{Connecting the LEDs together}
\begin{enumerate}
  \item Cut a ground wire long enough to connect the cathode of the central LED to the cathode of one peripheral LED. Strip both ends.
  \item Hold one end of the wire against the cathode of the central LED and solder it.
  \item Prepare a wire long enough to connect the peripheral LED of the lowest peripheral note to the adjacent LED: B to D in the large bowl, and C$\sharp$ to F$\sharp$ in the small bowl. Strip both ends.
  \item Hold both the wire already connected to the central note and the free wire against the cathode of the lowest peripheral note, B in the large bowl or C$\sharp$ in the small bowl, and solder them.
  \item Repeat from the previous step with each LED immediately to the left until all peripheral notes have been connected. Do not reconnect the final LED to the first one.
\end{enumerate}

\note{It is preferable to wait until the testing stage before shrinking the tubing around the solder joints.}

\section{Preparing the DB25 Cables}
\tutorial{https://youtu.be/j8yjBHBZ7oQ?is=TlLVSO2ycR7XFWDK}

\note{The tutorial uses a purpose-built DB25 cable tester designed to save time during group manufacturing workshops.}

\tools
\begin{itemize}
  \item Multimeter or cable tester
  \item Four-colour pen
\end{itemize}

\materials
\begin{itemize}
  \item Self-adhesive labels or writable adhesive tape
\end{itemize}

\procedure
\subsubsection{Preparing the identification labels}
\begin{itemize}
  \item Write the conductor numbers and note names on the labels using the following template: \url{https://drive.google.com/drive/folders/1wpoGsM1gzvmx8smcnHuNcGzMdusiLYeU}
  \item Separate the labels for each bowl to avoid incorrect distribution or confusion between notes.
\end{itemize}

\subsubsection{Cutting and stripping the cables}
\begin{itemize}
  \item Cut each DB25 cable 80 cm from the connector.
  \item Use a utility knife to remove 25 cm of the outer sheath, taking care not to cut the individual conductors.
  \item Strip every conductor.
\end{itemize}

\subsubsection{Identifying the conductors}
\begin{itemize}
  \item Select one conductor and hold it against one multimeter probe.
  \item With the other probe, test each pin of the DB25 connector until the matching pin is found.
  \item Determine the pin number using the wiring diagram at \url{https://docs.google.com/spreadsheets/d/1b5HFzodXUAxeUmf6LbNj5p8tK-nhCeLw/edit?usp=drivesdk&ouid=109005322539614617821&rtpof=true&sd=true}, then attach the corresponding label to the conductor.
  \item Repeat the operation, taking care not to mix the labels for the two bowls.
  \item Unidentified conductors in the B$\flat$ bowl cable may be shortened so that they do not interfere with the following stages.
\end{itemize}

\section{Connecting the DB25 Cables}
\tutorial{https://youtu.be/8Uv80K08HB4?is=GaTX8PpQ3HyFEN3Q}

\tools
\begin{itemize}
  \item Soldering iron
  \item Damp sponge
  \item Scissors
\end{itemize}

\materials
\begin{itemize}
  \item Solder
  \item Heat-shrink tubing
\end{itemize}

\procedure
\subsubsection{Preparation}
\begin{itemize}
  \item Tin all stripped conductors of the DB25 cable.
  \item Cut forty-two pieces of heat-shrink tubing, each approximately 1 cm long. When several tubing colours are available, prepare fourteen blue sleeves, fourteen red sleeves and fourteen sleeves in another chosen colour, referred to here as CC, which must be different from black, red and blue.
\end{itemize}

\subsubsection{Connecting the ground circuit to the DB25 cable}
\begin{itemize}
  \item Pass the black ground conductor of the DB25 cable through the heat-shrink sleeve on the LED that terminates the ground circuit: C in the large bowl and E$\flat$ in the small bowl.
  \item Align the end of the conductor with the existing solder joint and solder it in place.
\end{itemize}

\subsubsection{Connecting the LED and sensor conductors}
\begin{itemize}
  \item Select one conductor.
  \item Select a sleeve matching the colour used on its label. Use the CC sleeve for labels written in black, then pass the conductor through the sleeve.
  \item Following the colour code below, pass the LED lead or the sensor's red wire through the same sleeve, alongside the conductor and in the same direction.
  \item Align the conductor with the LED lead or sensor wire and solder the joint.
  \item Repeat this operation for every conductor.
\end{itemize}

\textbf{Colour code:}
\begin{itemize}
  \item Blue: short LED lead, positioned on the left
  \item Red: long LED lead, positioned on the right
  \item CC / black label text: red sensor wire
\end{itemize}

\section{Final Assembly}

\tools
\begin{itemize}
  \item Scissors
  \item Permanent marker suitable for smooth surfaces
  \item Drill
  \item 5 mm drill bit
  \item Hot-glue gun
  \item Heat source, such as a mini heat gun, lighter or soldering iron
\end{itemize}

\materials
\begin{itemize}
  \item Hot-glue stick
  \item Two plastic cable ties
  \item Camping mat or pieces of dense foam, 5--10 mm thick
  \item Double-sided adhesive tape
\end{itemize}

\procedure
\subsubsection{Positioning and shrinking the heat-shrink tubing}
\begin{itemize}
  \item After confirming that the device operates correctly by connecting it to an Arduino or ESP enclosure, position each sleeve so that it fully covers its solder joint.
  \item Apply heat briefly to shrink the tubing without melting the solder joints.
\end{itemize}

\subsubsection{Securing the DB25 cables to the bowls}
\begin{itemize}
  \item Position the DB25 cable between the two lowest peripheral notes: B and C in the large bowl, or C$\sharp$ and E$\flat$ in the small bowl.
  \item Mark two points, one on each side of the cable.
  \item Move the cable aside and drill at the marked points.
  \item Remove the dust with a cloth or vacuum cleaner.
  \item Feed the narrow end of the cable tie through one hole from the inside, then bring the same end back through the second hole from the outside.
  \item Position the DB25 cable between the two holes.
  \item Close and tighten the cable tie.
  \item Adjust the cable length so that approximately 5 cm of unstripped cable extends beyond the cable tie towards the centre of the bowl.
  \item Apply hot glue and allow it to set.
\end{itemize}

\subsubsection{Making and attaching the feet}
\begin{itemize}
  \item Cut eight pieces measuring 5 cm by 1.5 cm from the camping mat or dense foam.
  \item For each foot, apply double-sided adhesive tape or hot glue to one face. Remove the protective film when using adhesive tape.
  \item Attach the foot to the rim of the bowl, avoiding the DB25 cable.
  \item Repeat until four feet have been evenly distributed around each bowl.
\end{itemize}

\end{document}
